\documentclass[11pt,a4paper]{article}
\usepackage[margin=22mm]{geometry}
\usepackage{amsmath,amssymb}
\usepackage{fontspec}
\usepackage{polyglossia}
\setdefaultlanguage{russian}
\setmainfont{DejaVu Serif}
\usepackage[unicode,hidelinks]{hyperref}

\allowdisplaybreaks[1]
\raggedright
\setlength{\parskip}{0.35em}
\newcommand{\dd}{\mathop{}\!\mathrm{d}}
\newcommand{\avg}[1]{\langle #1\rangle}
\newcommand{\ket}[1]{\lvert #1\rangle}
\newcommand{\bra}[1]{\langle #1\rvert}
\newcommand{\braket}[2]{\langle #1\mid #2\rangle}
\DeclareMathOperator{\Var}{Var}

\newcommand{\solutiontitle}[4]{%
  \hypersetup{pdftitle={Квантовая механика 1. Решения ДЗ №#1}}%
  \begin{center}
    {\small Квантовая механика 1 \enspace\textperiodcentered\enspace Семестр 1\par}
    \vspace{1ex}
    {\LARGE\bfseries Решения домашнего задания №#1\par}
    \vspace{0.5ex}
    {\large #2\par}
  \end{center}
  \noindent{\small Условия приведены в файле «#3», #4.}\par\medskip
}

\begin{document}
\solutiontitle{5}{Импульсные волновые функции}{ДЗ2026(1)\_05.pdf}{с.~1--2}

Ядро прямого преобразования в \(d\) измерениях равно
\((2\pi\hbar)^{-d/2}e^{-i\mathbf p\cdot\mathbf r/\hbar}\).
В обратном меняется знак показателя. В переменной
\(\mathbf k=\mathbf p/\hbar\) пишем симметричную пару с
\((2\pi)^{-d/2}\) и \(\exp(\mp i\mathbf k\cdot\mathbf r)\).

\section{Сдвинутая гауссова функция}

Нормировка даёт \(1=C^2\sigma\sqrt\pi\), \(\sigma>0\).
Сдвиг \(y=x-x_0\) и гауссов интеграл
\(\int e^{-y^2/(2\sigma^2)-iky}\dd y
=\sqrt{2\pi}\sigma e^{-\sigma^2 k^2/2}\) дают
\[
 C=(\sigma\sqrt\pi)^{-1/2},\qquad
 a(p)=\frac{C\sigma}{\sqrt\hbar}
       e^{-ipx_0/\hbar}e^{-\sigma^2p^2/(2\hbar^2)}.
\]
В координатном представлении \(\Psi'=-y\Psi/\sigma^2\), поэтому
интеграл \(-i\hbar\int\Psi^*\Psi'\dd x\) обнуляется по нечётности
по \(y\). Интегрируя по частям, находим
\(\avg{p^2}=\hbar^2\int|\Psi'|^2\dd x
=\hbar^2\avg{y^2}/\sigma^4=\hbar^2/(2\sigma^2)\).
Здесь \(\avg{y^2}=\sigma^2/2\) получается интегрированием
гауссовой плотности по частям.
В импульсном представлении \(|a(p)|^2\) чётна, а замена
\(u=\sigma p/\hbar\) даёт ту же нормированную плотность
\(e^{-u^2}/\sqrt\pi\) с \(\avg{u^2}=1/2\). В обоих представлениях
\(\avg{p}=0,\ \Var(p)=\hbar^2/(2\sigma^2)\).

\section{Экспонента с изломом в начале координат}

При \(a>0\) норма равна \(2C^2\int_0^\infty e^{-2x/a}\dd x=aC^2\).
Разбив прямое преобразование на полуоси, получаем
\begin{align*}
 C&=a^{-1/2},\\
 a(p)&=\frac{C}{\sqrt{2\pi\hbar}}
   \left(\frac1{a^{-1}+ip/\hbar}+\frac1{a^{-1}-ip/\hbar}\right)
   =\sqrt{\frac{2a}{\pi\hbar}}\frac1{1+(ap/\hbar)^2}.
\end{align*}
Средний импульс нулевой по чётности плотности.
Для \(u=ap/\hbar\) интеграл
\(\int_\mathbb R u^2(1+u^2)^{-2}\dd u=\pi/2\) следует из
\((1+u^2)^{-1}-(1+u^2)^{-2}\) и подстановки \(u=\tan v\).
Значит \(\avg{p^2}=(2/\pi)(\hbar/a)^2(\pi/2)=\hbar^2/a^2\).
В координатах слабая производная
\(\Psi'=-(\operatorname{sgn}x)\Psi/a\) лежит в \(L^2\).
По Планшерелю второй момент равен
\(\hbar^2\|\Psi'\|^2=\hbar^2/a^2\), и первый момент равен нулю
по нечётности \(\Psi^*\Psi'\). Второй момент здесь вычисляется
как квадратичная форма, поскольку у \(\Psi''\) есть дельта-вклад.
Дисперсия \(\Var(p)=\hbar^2/a^2\).

\section{Нечётная импульсная функция с гауссовой огибающей}

В \(a(p)=Cp e^{-p^2/(2p_0^2)}\) берём \(p_0>0\), либо заменяем
\(p_0\) на \(|p_0|\).
Из \(\int p^2e^{-p^2/p_0^2}\dd p=p_0^3\sqrt\pi/2\)
следует \(C=\sqrt2/(\pi^{1/4}p_0^{3/2})\).
Обозначим \(\ell=\hbar/p_0\). В обратном преобразовании
множитель \(p\) заменяется на \(-i\hbar\partial_x\) после
преобразования гауссианы. Поэтому
\[
 \Psi(x)=-i\hbar C\frac{p_0}{\sqrt\hbar}\frac{\dd}{\dd x}
               e^{-x^2/(2\ell^2)}
        =\frac{i\sqrt2}{\pi^{1/4}\ell^{3/2}}\,x e^{-x^2/(2\ell^2)}.
\]
Плотности \(|a(p)|^2\) и \(|\Psi(x)|^2\) чётны.
Для любого масштаба \(s>0\) отношение
\(\int v^4 e^{-v^2/s^2}\dd v/\int v^2 e^{-v^2/s^2}\dd v
=3s^2/2\) получается одним интегрированием по частям. Следовательно
\[
 \avg{x}=\avg{p}=0,\qquad
 \Var(x)=\frac{3\hbar^2}{2p_0^2},\qquad \Var(p)=\frac{3p_0^2}{2}.
\]

\section{Трёхмерные обобщённые состояния}

\textbf{(a)} В обратном преобразовании \(\sin(k_xa)\) даёт
разность дельта-функций по \(x\), \(e^{-ik_y b}\) фиксирует \(y=b\),
а \(\delta(k_z-k_0)\) даёт плоскую волну по \(z\). По интегралам
\(\int e^{iku}\dd k=2\pi\delta(u)\) имеем
\[
 \Psi(\mathbf r)=\frac{C\sqrt{2\pi}}{2i}
       \bigl[\delta(x+a)-\delta(x-a)\bigr]\delta(y-b)e^{ik_0z}.
\]
При \(a\ne0\) ненулевые компоненты поддержаны на двух разных
значениях \(x\). Множитель \(e^{ik_0z}\) имеет определённый
импульс \(p_z=\hbar k_0\). В импульсном представлении \(k_x,k_y\)
меняются непрерывно. По \(z\) плоская волна имеет постоянный модуль.
Таким образом определены только \(y=b\) и \(p_z=\hbar k_0\).
Дельта-функции и плоская волна исключают нормировку в \(L^2\).
В смешанном дельта-нормированном базисе два слагаемых с \(x=\pm a\)
имеют равные по модулю относительные амплитуды. Например, замена
обеих координатных дельта-функций прямоугольниками площади один
ширины \(\varepsilon_x,\varepsilon_y\) и ограничение плоской волны
длиной \(L_z\) дают нормированное приближение, если умножить
записанную разность, прямоугольник по \(y\) и волну по \(z\)
на \(\sqrt{\varepsilon_x\varepsilon_y/(2L_z)}\).
При \(a=0\) исходная функция нулевая.

\textbf{(b)} Пусть \(x_0>0\), а \(k\) в \(e^{-iky}\cos(kz)\)
есть фиксированное число. Разделяем интегралы по координатам.
\begin{align*}
 \int e^{-|x|/x_0-ik_xx}\dd x&=\frac{2x_0}{1+x_0^2k_x^2},\\
 \int e^{-i(k_y+k)y}\dd y&=2\pi\delta(k_y+k),\\
 \int e^{-ik_zz}\cos(kz)\dd z&=\pi[\delta(k_z-k)+\delta(k_z+k)].
\end{align*}
Произведение с мерой \((2\pi)^{-3/2}\) даёт
\[
 a(\mathbf k)=\frac{C\sqrt{2\pi}\,x_0}{1+x_0^2 k_x^2}
             \delta(k_y+k)\,[\delta(k_z-k)+\delta(k_z+k)].
\]
При \(k\ne0\) единственная заданная величина из списка
есть \(p_y=-\hbar k\). По \(z\) присутствуют оба импульса
\(\pm\hbar k\) с равными амплитудами, по \(x\) спектр непрерывен,
а ни одна координата не фиксирована. При \(k=0\) также \(p_z=0\).
Норма на \(\mathbb R^3\) расходится. В ящике по \(y,z\)
длин \(L_y,L_z\) с периодами, совместимыми с \(k\),
\(\int e^{-2|x|/x_0}\dd x=x_0\) и
\(\int_0^{L_z}\cos^2(kz)\dd z=L_z/2\) при \(k\ne0\).
Отсюда \(|C_{\rm box}|=\sqrt{2/(x_0L_yL_z)}\), а при \(k=0\)
получаем \(|C_{\rm box}|=(x_0L_yL_z)^{-1/2}\).

\section{Сферически-симметричное состояние}

При \(a>0\) нормировка \(4\pi C^2\int_0^\infty r^2e^{-2r/a}\dd r=1\)
даёт \(C=(\pi a^3)^{-1/2}\). Выбираем полярную ось вдоль
\(\mathbf p\), где \(p=|\mathbf p|\). Угловой интеграл равен
\(4\pi\sin(kr)/(kr)\), \(k=p/\hbar\).
Из \(\int_0^\infty e^{-br}\sin(kr)\dd r=k/(b^2+k^2)\)
дифференцированием по \(b=1/a\) находим
\[
 a(\mathbf p)=\frac{4\pi C}{(2\pi\hbar)^{3/2}}
                \frac{2/a}{(a^{-2}+k^2)^2}
       =\frac{2\sqrt2}{\pi}\left(\frac a\hbar\right)^{3/2}
                \frac1{[1+(ap/\hbar)^2]^2}.
\]
Симметрия даёт \(\avg{\mathbf p}=\mathbf0\). По Планшерелю
\(\avg{\mathbf p^2}=\hbar^2\int|\nabla\Psi|^2\dd^3r
=\hbar^2/a^2\), так как \(\nabla\Psi=-\Psi\,\mathbf r/(ar)\)
почти всюду. Поэтому дисперсия каждой компоненты равна
\(\hbar^2/(3a^2)\), а сумма трёх дисперсий равна \(\hbar^2/a^2\).

Для модуля \(p\) мера углового интегрирования даёт
плотность \(32u^2/[\pi(1+u^2)^4]\) по \(u=ap/\hbar\).
Подстановка \(v=1+u^2\) даёт
\(\int_0^\infty u^3(1+u^2)^{-4}\dd u=1/12\).
С учётом того же второго момента \(\avg{p^2}\)
получаем
\[
 \avg{p}=\frac{8\hbar}{3\pi a},\qquad
 \Var(p)=\frac{\hbar^2}{a^2}\left(1-\frac{64}{9\pi^2}\right).
\]

\section{Чётность и трансляция (6*)}

В прямом интеграле для \(Pf(x)=f(-x)\) замена \(u=-x\)
меняет знак в ядре, значит \((Pa)(p)=a(-p)\).
Для \(T_\delta f(x)=f(x+\delta)\) подстановка \(u=x+\delta\)
порождает фазу \(e^{ip\delta/\hbar}\). Поэтому
\[
 (Pa)(p)=a(-p),\qquad (T_\delta a)(p)=e^{ip\delta/\hbar}a(p).
\]

\end{document}
